\setlength{\parindent}{0pt}
\small

\noindent
\begin{minipage}[t]{0.485\textwidth}
  \vspace{0pt}
  \setlength{\parskip}{3.5pt}

  {\color{stewartcyan}\textbf{15--18}}\enspace Kiểm tra xem hàm số có thỏa mãn các giả thiết của Định lý giá trị trung bình trên đoạn đã cho hay không. Sau đó tìm tất cả các số $c$ thỏa mãn kết luận của Định lý giá trị trung bình.

  \textbf{15.}\enspace $f(x) = 2x^2 - 3x + 1, \quad [0, 2]$

  \textbf{16.}\enspace $f(x) = x^3 - 3x + 2, \quad [-2, 2]$

  \noindent
  \textbf{17.}\enspace $f(x) = \ln x, \quad [1, 4]$
  \hfill
  \textbf{18.}\enspace $f(x) = 1/x, \quad [1, 3]$

  \exerciseline

  \graphingicon\enspace {\color{stewartcyan}\textbf{19--20}}\enspace Tìm số $c$ thỏa mãn kết luận của Định lý giá trị trung bình trên khoảng đã cho. Vẽ đồ thị của hàm số, đường cát tuyến đi qua các điểm mút và tiếp tuyến tại $(c, f(c))$. Cát tuyến và tiếp tuyến có song song với nhau không?

  \noindent
  \textbf{19.}\enspace $f(x) = \sqrt{x}, \quad [0, 4]$
  \hfill
  \textbf{20.}\enspace $f(x) = e^{-2x}, \quad [0, 2]$

  \exerciseline

  \textbf{21.}\enspace Cho $f(x) = (x - 3)^{-2}$. Chứng minh rằng không có giá trị nào của $c$ trong khoảng $(1, 4)$ sao cho $f(4) - f(1) = f'(c)(4 - 1)$. Tại sao điều này không mâu thuẫn với Định lý giá trị trung bình?

  \textbf{22.}\enspace Cho $f(x) = 2 - |2x - 1|$. Chứng minh rằng không có giá trị nào của $c$ sao cho $f(3) - f(0) = f'(c)(3 - 0)$. Tại sao điều này không mâu thuẫn với Định lý giá trị trung bình?

  {\color{stewartcyan}\textbf{23--24}}\enspace Chứng minh rằng phương trình có đúng một nghiệm thực.

  \noindent
  \textbf{23.}\enspace $2x + \cos x = 0$
  \hfill
  \textbf{24.}\enspace $x^3 + e^x = 0$

  \exerciseline

  \textbf{25.}\enspace Chứng minh rằng phương trình $x^3 - 15x + c = 0$ có nhiều nhất một nghiệm trong đoạn $[-2, 2]$.

  \textbf{26.}\enspace Chứng minh rằng phương trình $x^4 + 4x + c = 0$ có nhiều nhất hai nghiệm thực.

  \textbf{27.}\enspace \textbf{(a)}~Chứng minh rằng một đa thức bậc 3 có nhiều nhất ba nghiệm thực.\par
  \textbf{(b)}~Chứng minh rằng một đa thức bậc $n$ có nhiều nhất $n$ nghiệm thực.

  \textbf{28.}\enspace \textbf{(a)}~Giả sử $f$ khả vi trên $\mathbb{R}$ và có hai nghiệm. Chứng minh rằng $f'$ có ít nhất một nghiệm.\par
  \textbf{(b)}~Giả sử $f$ khả vi hai lần trên $\mathbb{R}$ và có ba nghiệm. Chứng minh rằng $f''$ có ít nhất một nghiệm thực.\par
  \textbf{(c)}~Bạn có thể tổng quát hóa các phần (a) và (b) không?

  \textbf{29.}\enspace Nếu $f(1) = 10$ và $f'(x) \ge 2$ với $1 \le x \le 4$, thì giá trị nhỏ nhất mà $f(4)$ có thể nhận là bao nhiêu?
\end{minipage}%
\hfill
\begin{minipage}[t]{0.485\textwidth}
  \vspace{0pt}
  \setlength{\parskip}{3.5pt}

  \textbf{30.}\enspace Giả sử rằng $3 \le f'(x) \le 5$ với mọi giá trị của $x$. Chứng minh rằng $18 \le f(8) - f(2) \le 30$.

  \textbf{31.}\enspace Có tồn tại một hàm số $f$ sao cho $f(0) = -1$, $f(2) = 4$, và $f'(x) \le 2$ với mọi $x$ không?

  \textbf{32.}\enspace Giả sử $f$ và $g$ liên tục trên $[a, b]$ và khả vi trên $(a, b)$. Giả sử thêm rằng $f(a) = g(a)$ và $f'(x) < g'(x)$ với $a < x < b$. Chứng minh rằng $f(b) < g(b)$.\enspace [\textit{Gợi ý:} Áp dụng Định lý giá trị trung bình cho hàm số $h = f - g$.]

  \textbf{33.}\enspace Chứng minh rằng $\sin x < x$ nếu $0 < x < 2\pi$.

  \textbf{34.}\enspace Giả sử $f$ là hàm số lẻ và khả vi ở mọi nơi. Chứng minh rằng với mọi số dương $b$, tồn tại một số $c$ thuộc $(-b, b)$ sao cho $f'(c) = f(b)/b$.

  \textbf{35.}\enspace Sử dụng Định lý giá trị trung bình để chứng minh bất đẳng thức
  \[
    |\sin a - \sin b| \le |a - b| \qquad \text{với mọi } a \text{ và } b
  \]

  \textbf{36.}\enspace Nếu $f'(x) = c$ ($c$ là hằng số) với mọi $x$, hãy dùng Hệ quả 7 để chứng minh rằng $f(x) = cx + d$ với một hằng số $d$ nào đó.

  \textbf{37.}\enspace Cho $f(x) = 1/x$ và
  \[
    g(x) = \begin{cases}
      \dfrac{1}{x} & \text{nếu } x > 0 \\[6pt]
      1 + \dfrac{1}{x} & \text{nếu } x < 0
    \end{cases}
  \]
  Chứng minh rằng $f'(x) = g'(x)$ với mọi $x$ thuộc tập xác định của chúng. Ta có thể kết luận từ Hệ quả 7 rằng $f - g$ là một hàm hằng không?

  {\color{stewartcyan}\textbf{38--39}}\enspace Sử dụng phương pháp của Ví dụ 6 để chứng minh đồng nhất thức.

  \textbf{38.}\enspace $\arctan x + \arctan\left(\dfrac{1}{x}\right) = \dfrac{\pi}{2}, \quad x > 0$

  \textbf{39.}\enspace $2\sin^{-1}x = \cos^{-1}(1 - 2x^2), \quad x \ge 0$

  \textbf{40.}\enspace Lúc 2:00 chiều, đồng hồ tốc độ của một xe ô tô chỉ $30\text{ mi/h}$. Lúc 2:10 chiều, nó chỉ $50\text{ mi/h}$. Chứng minh rằng tại một thời điểm nào đó giữa 2:00 và 2:10, gia tốc của xe đúng bằng $120\text{ mi/h}^2$.

  \textbf{41.}\enspace Hai vận động viên bắt đầu cuộc đua cùng một lúc và về đích cùng nhau (hòa). Chứng minh rằng tại một thời điểm nào đó trong cuộc đua, họ có cùng tốc độ.\enspace [\textit{Gợi ý:} Xét hàm số $f(t) = g(t) - h(t)$, trong đó $g$ và $h$ là hàm vị trí của hai vận động viên.]

  \textbf{42.}\enspace \textbf{Điểm bất động}\enspace Một số $a$ được gọi là một \textit{điểm bất động} của hàm số $f$ nếu $f(a) = a$. Chứng minh rằng nếu $f'(x) \ne 1$ với mọi số thực $x$, thì $f$ có nhiều nhất một điểm bất động.
\end{minipage}

\vspace{22pt}

\noindent
\begin{minipage}[t]{0.22\textwidth}
  \vspace{0pt}
  \begin{tikzpicture}[baseline=8pt]
    \foreach \y in {0pt, 3.2pt, 6.4pt, 9.6pt} {
      \draw[color=stewartcyan, line width=1.2pt] (0, \y) -- (32pt, \y);
    }
    \node[anchor=west, inner sep=0pt] at (38pt, 4.8pt) {%
      \fontsize{24}{26}\selectfont\sffamily\bfseries\color{stewartcyan} 4.3%
    };
  \end{tikzpicture}
\end{minipage}%
\hfill
\begin{minipage}[t]{0.75\textwidth}
  \vspace{0pt}
  {\fontsize{14.5}{18}\selectfont\sffamily\bfseries Đạo hàm cho ta biết điều gì về hình dạng của đồ thị}\par
  \vspace{8pt}
  \normalsize
  Nhiều ứng dụng của giải tích phụ thuộc vào khả năng của chúng ta trong việc suy ra các tính chất về một hàm số $f$ từ những thông tin liên quan đến các đạo hàm của nó. Bởi vì $f'(x)$ biểu thị hệ số góc của đường cong $y = f(x)$ tại điểm $(x, f(x))$, nó cho chúng ta biết hướng đi mà đường cong
\end{minipage}
